\documentclass{article}
\usepackage{../Self_Style}

\usepackage[normalem]{ulem} %for strike through

\title{Dualizability in Symmetric Monoidal Category}
\author{Zih-Yu Hsieh}
\date{\today}

\begin{document}
\maketitle
%\tableofcontents

This is devoted for the categorical details in 

We'll follow the notations Aditya covered in the talk. Fix $(\mathcal{C},\tensor,I)$ as our symmetric monoidal category.

\begin{defn}
    Given an object $A\in\mathcal{C}$, an object $B\in\mathcal{C}$ is called $A$'s \emph{weak dual}, if there exists a natural isomorphism $\xi:\mathcal{C}(\_\tensor A,I)\xrightarrow{\sim} \mathcal{C}(\_,B)$ between functors $\mathcal{C}^\Op\rightarrow \Sets$. In this case, $A$ is said to be \emph{weakly dualizable}.
\end{defn}

\hfil

One important fact one needs to verify is the uniqueness (up to isomorphism), this is indeed true:

\begin{prop}
    If $A\in\mathcal{C}$ is weakly dualizable, its weak dual is unique up to isomorphism.
\end{prop}
\begin{proof}
    This is a Yoneda Lemma type argument. Let $B,B'\in\mathcal{C}$ be two weak duals of $A$, then there is a natural isomorphism $\xi:\mathcal{C}(\_,B)\xrightarrow{\sim}\mathcal{C}(\_\tensor A,I)\xrightarrow{\sim}\mathcal{C}(\_,B')$ (abuse of notation, still use $\xi$ as natural isomorphism).

    Using the isomorphisms $\xi_B:\mathcal{C}(B,B)\xrightarrow{\sim}\mathcal{C}(B,B')$ and $\xi_{B'}^{-1}:\mathcal{C}(B',B')\xrightarrow{\sim}\mathcal{C}(B',B)$, we get the following two morphisms:
    \[\eta:= \xi_B(\id_B):B\rightarrow B',\quad \mu:=\xi_{B'}^{-1}(\id_{B'}):B'\rightarrow B\]
    Then, consider the following two diagrams using the naturality:
    \[\xymatrix{
        \mathcal{C}(B,B)\ar[r]^-{\xi_B}_-{\sim} \ar[d]_-{\_\circ\mu} & \mathcal{C}(B,B') \ar[d]^-{\_\circ\mu}\\
        \mathcal{C}(B',B')\ar[r]^-{\sim}_-{\xi_{B'}} & \mathcal{C}(B',B')
    },\quad\quad \xymatrix{
        \id_B \ar@{|->}[r] \ar@{|->}[d] & \eta \ar@{|->}[d]\\
        \mu \ar@{|->}[r] & \id_{B'}=\eta\circ \mu
    }\]
    \[\xymatrix{
        \mathcal{C}(B',B')\ar[r]^-{\xi_B^{-1}}_-{\sim} \ar[d]_-{\_\circ\eta} & \mathcal{C}(B',B) \ar[d]^-{\_\circ\eta}\\
        \mathcal{C}(B,B')\ar[r]^-{\sim}_-{\xi_{B'}^{-1}} & \mathcal{C}(B,B)
    },\quad\quad \xymatrix{
        \id_{B'} \ar@{|->}[r] \ar@{|->}[d] & \mu \ar@{|->}[d]\\
        \eta \ar@{|->}[r] & \id_{B}=\mu\circ \eta
    }\]
    This shows $\mu,\eta$ are mutual inverses, hence $B\cong B'$.
\end{proof}

If the weakly dualizable objects are small, then there's no ambiguity choosing a representative of $A$, denote as $DA$.

\newpage

Another notion is how to ``formally evaluate'' weakly dualizable objects:
\begin{defn}
    Given $A\in\mathcal{C}$ that's weakly dualizable, consider the isomorphism 
    \[\xi_{DA}:\mathcal{C}(DA\tensor A,I)\xrightarrow{\sim}\mathcal{C}(DA,DA)\]
    Define the \emph{evaluation morphism} $\epsilon_A:DA\tensor A\rightarrow I$ as the dual to the following identity map, $\epsilon_A:= \xi_{DA}^{-1}(\id_{DA})$.
\end{defn}

\hfil

Now, let $\mathcal{C}^*$ denote the full subcategory of all weakly dualizable objects, the aim is to show dualization is functorial, i.e., there is a functor $D:(\mathcal{C}^*)^\Op\rightarrow\Sets$. It's clear that $DA$ is well-defined (up to isomorphism), so the morphism (and preservation of composition) is the actual problem.

\begin{defn}
    Given a morphism $f:A\rightarrow A'$ in $\mathcal{C}^*$, consider the morphism 
    \[DA'\tensor A\xrightarrow{\id_{DA'}\tensor f}DA'\tensor A'\xrightarrow{\epsilon_{A'}}I\]
    Using the isomorphism $\xi_{DA'}:\mathcal{C}(DA'\tensor A,I)\xrightarrow{\sim}\mathcal{C}(DA',DA)$, define 
    \[Df:= \xi_{DA'}(\epsilon_{A'}\circ (\id_{DA'}\tensor f)):DA'\rightarrow DA\]
\end{defn}

\hfil

The goal is to show this dualization of morphism is functorial. Before that, fix $f:A\rightarrow A'$ in $\mathcal{C}^*$, and $Df:DA'\rightarrow DA$ in $\mathcal{C}$, consider the following commutative diagram using the natural isomorphism $\xi:\mathcal{C}(\_\tensor A,I)\xrightarrow{\sim}\mathcal{C}(\_,DA)$:
\[\xymatrix{
    \mathcal{C}(DA\tensor A,I) \ar[d]_-{\_\circ (Df\tensor \id_A)} \ar[r]^-{\xi_{DA}}_-{\sim} & \mathcal{C}(DA,DA) \ar[d]^-{\_\circ Df}\\
    \mathcal{C}(DA'\tensor A,I) \ar[r]^-{\sim}_-{\xi_{DA'}} & \mathcal{C}(DA',DA)
},\quad\quad \xymatrix{
        \epsilon_A \ar@{|->}[r] \ar@{|->}[d] & \id_{DA} \ar@{|->}[d]\\
        \epsilon_A\circ (Df\tensor \id_A) \ar@{|->}[r] & Df
    }\]
So, we see the equality $\xi_{DA'}(\epsilon_A\circ (Df\tensor \id_A)) = Df$, and don't forget that $Df:= \xi_{DA'}(\epsilon_{A'}\circ(\id_{DA'}\tensor f))$. So, the isomorphism deduces $\epsilon_A\circ (Df\tensor \id_A) = \epsilon_{A'}\circ (\id_{DA'}\tensor f)$. 

Using this, we deduce the funcoriality as morphism $DA''\tensor A\rightarrow I$:
\begin{lem}
    $D:(\mathcal{C}^*)^\Op\rightarrow\Sets$ is functorial, namely $D(g\circ f)= Df\circ Dg$.
\end{lem}
\begin{proof}
    Given $f:A\rightarrow A'$ and $g:A'\rightarrow A''$ in $\mathcal{C}^*$, we get the following equality:
    \[\epsilon_{A''}\circ (\id_{DA''}\tensor (g\circ f)) = \epsilon_{A''}\circ (\id_{DA''}\tensor g)\circ (\id_{DA''}\tensor f)= \epsilon_{A'}\circ (Dg\tensor \id_{A'})\circ (\id_{DA''}\tensor f)\]
    \[= \epsilon_{A'}\circ (Dg\tensor f) = \epsilon_{A'}\circ (\id_{DA'}\tensor f)\circ (Dg\tensor \id_A) \]
    \[= \epsilon_A\circ (Df\tensor \id_A)\circ (Dg\tensor \id_A)= \epsilon_A\circ ((Df\circ Dg)\tensor \id_A)\]
    Hence, apply $\xi_{DA''}$, we get:
    \[Df\circ Dg = \xi_{DA''}(\epsilon_A\circ (Df\circ Dg)\tensor \id_A) = \xi_{DA''}(\epsilon_{A''}\circ (\id_{DA''}\tensor (g\circ f)))= D(g\circ f)\]
    If one is interested they can verify the above are all well-defined :)
\end{proof}


\end{document}